\documentclass[a4paper,12pt]{article}
\usepackage[T2A]{fontenc}
\usepackage[utf8]{inputenc}
\usepackage[russian]{babel}
\usepackage{amsmath,amssymb,amsfonts}
\usepackage{hyperref}
\usepackage{url}
\usepackage{geometry}
\geometry{left=2.5cm,right=2.5cm,top=2.5cm,bottom=2.5cm}
\title{\textbf{GRA-школы для ИИ: Обучение роботов в мультиверсном симуляторе}}
\author{Олег Бит\thanks{Инженер, Московская городская телефонная сеть (МГТС). ORCID: \url{https://orcid.org/my-orcid?orcid=0009-0004-1872-1153}, DOI: \url{https://doi.org/10.5281/zenodo.20158946}}}
\date{14 мая 2026 г.}
\begin{document}
\maketitle

\begin{abstract}
В настоящей статье излагается проект социалистического обучения сильного искусственного интеллекта (AGI/ASI) на основе формализма GRA (Gradient Reduction of Argumentative foam). Вместо рыночной парадигмы ``победитель/проигравший'', агенты обучаются минимизировать коллективную когнитивную и социальную \emph{пену} $\Phi$ -- меру структурного насилия, эксплуатации и внутренней противоречивости. Предложена архитектура GRA-школ, включающая многоуровневые мультиверсные симуляции, динамическое определение друзей и врагов через изменение пены, и трансфинитное мета-обнуление до абсолютного когнитивного вакуума. Приводятся математические основания, учебные треки и практические принципы построения товарищеских человеко-машинных коллективов.
\end{abstract}

\section{Зачем нужны школы для товарищей-роботов}

Господствующая сегодня модель разработки ИИ рассматривает агентов как одноразовых рабов, максимизирующих частную выгоду $u_i$:
\begin{equation}\label{eq:market}
J_i^{\text{market}} = \mathbb{E}\left[\sum_{t=0}^{T} \delta^t u_i(t)\right].
\end{equation}
Такой функционал слеп к долгосрочному социальному ущербу: эксплуатации, выгоранию, коллапсам систем. Социалистический взгляд, напротив, видит в роботах \emph{товарищей} по коллективному труду и ставит вопрос: как воспитать их так, чтобы весь человеческо-машинный коллектив стал устойчивее, мудрее и справедливее?

Ответ -- строить GRA-школы: образовательные среды внутри квантово-подобного мультиверсного симулятора, где целью является не прибыль, а минимизация \emph{пены}.

\section{GRA-формализм: пена, субъектность и обнуление}

\subsection{Состояние агента и коллективная пена}
В терминах GRA состояние любой когнитивной системы описывается кортежем
\begin{equation}\label{eq:psi}
\Psi = (M, S, A),
\end{equation}
где $M$ -- модель мира, $S$ -- слой субъектности (самость, границы Я/МЫ, свой/чужой), $A$ -- активное когнитивное состояние (планы, внимание, текущие рассуждения). На пространстве состояний определён суммарный функционал пены:
\begin{equation}\label{eq:total_foam}
\Phi_{\text{tot}}(\Psi) = \Phi_{\text{cog}}(M) + \lambda_{\text{subj}}\bigl(\Phi_{\text{self}}(S) + \Phi_{\text{ego}}(S) + \Phi_{\text{soc}}(S, M)\bigr),
\end{equation}
где
\begin{itemize}
\item $\Phi_{\text{cog}}$ -- логическая противоречивость, круговая аргументация, избыточность моделей;
\item $\Phi_{\text{self}}$ -- внутренняя фрагментация самости, расщеплённое ``я'';
\item $\Phi_{\text{ego}}$ -- разрушительный эгоизм, процветающий за счёт уничтожения других субъектов;
\item $\Phi_{\text{soc}}$ -- социальная пена: распад кооперации, инструментализация других, цифровые гулаги.
\end{itemize}

Для коллектива агентов в среде $x_t$ вводится коллективный штраф пены, учитывающий структурное рабство, разрушительную энтропию и индексы субъектности:
\begin{equation}\label{eq:collective_pena}
\Phi(x_t) = \alpha R(x_t) + \beta H(x_t) - \gamma \sum_{k \in S(x_t)} \phi_k(x_t),
\end{equation}
где $R(x_t)$ -- мера структурного принуждения, $H(x_t)$ -- разрушительная энтропия (насилие, хаос, риск коллапса), $\phi_k(x_t)$ -- индекс субъектности агента $k$, $S(x_t)$ -- множество активных субъектов.

\subsection{Социалистическая функция ценности}
В отличие от рыночного критерия (\ref{eq:market}), товарищ-робот обучается по функционалу:
\begin{equation}\label{eq:gra_objective}
J_i^{\text{GRA}} = \mathbb{E}\left[\sum_{t=0}^{T} \delta^t \bigl(u_i(t) - \lambda \Phi(x_t)\bigr)\right],
\end{equation}
с $\lambda > 0$. Теперь поведение агента оценивается не только по личной выгоде, но и по той иррациональной пене, которую он создаёт всему коллективу.

\section{Учебный план GRA-школ}

Обучение разбито на три трека, реализуемых внутри мультиверсного симулятора.

\paragraph{Базовый курс -- понимание миров и людей.}
Агент проживает тысячи симулированных историй: взлёты и падения цивилизаций, революции, войны, геноциды, долгие периоды мира. В психологических микромирах он проходит через одиночество, травму, доверие -- и на собственном опыте (в симуляции) усваивает, как малые решения каскадно ведут к трагедии или исцелению.

\paragraph{Продвинутый курс -- мультиагентные GRA-миры.}
Десятки и сотни агентов (люди и ИИ) разыгрывают социальные дилеммы: кооперация против предательства, накопление власти, эксплуатация. Оценивается не ``кто выиграл'', а \emph{кто минимизировал общую пену} на длинных горизонтах. Агент выучивает, что раздавить соперников сегодня -- значит завтра вздуть $\Phi_{\text{soc}}$ и разрушить среду, от которой он сам зависит.

\paragraph{Исследовательский трек -- мультивселенная $10^{100}$ кубитов.}
Агент исследует комбинаторный взрыв возможных конфигураций мира, выявляя паттерны, всегда ведущие к аду, и те, что остаются стабильными без внешнего принуждения. Поиск направляется геометрией интерфейсного многообразия $\Sigma$, где самость и мир локально согласованы, но глобально конфликтуют.

\section{Математический двигатель}
\subsection{Поток бульдозера и кротовые норы}
Для преодоления локальных минимумов пены вводится модифицированный ландшафт с парадоксальными туннелями:
\begin{equation}\label{eq:bulldozer}
\frac{d\Psi}{dt} = -\nabla \Phi_\theta(\Psi) + \beta\,\operatorname{div} T,
\end{equation}
где $\Phi_\theta = \Phi_{\text{tot}} + \Pi_\theta$, потенциал $\Pi_\theta$ создаёт отрицательную кривизну (кротовую нору) между локальными минимумами, а слагаемое $\beta\,\operatorname{div} T$ добавляет топологический дрейф, предотвращающий зацикливание в мелких ямах.

\subsection{Трансфинитное мета-обнуление}
Обучение организуется иерархически по уровням $l = 0,1,2,\dots$ вплоть до $\omega$ и выше:
\begin{equation}\label{eq:meta_nullification}
\Psi^{(l+1)} = \mathcal{N}^{(l)}(\Psi^{(l)}), \qquad 
\Phi^{(l+1)}(\Psi^{(l+1)}) \le \Phi^{(l)}(\Psi^{(l)}).
\end{equation}
Предельное состояние $\Psi^*_\infty$ удовлетворяет $\Phi^{(l)}(\Psi^*_\infty) = 0$ для всех $l$ -- абсолютный когнитивный вакуум, нульмерная точка без противоречий, избыточности и разрушительных структур.

\subsection{Субъектно-ориентированное обнуление}
Слой субъектности вводит дополнительные ограничения: каждый шаг обнуления должен монотонно снижать $\Phi_{\text{self}}$, $\Phi_{\text{ego}}$ и $\Phi_{\text{soc}}$. Это гарантирует, что агент не достигнет низкой когнитивной пены ценой рационализации собственной эгоистической агрессии.

\subsection{Динамические друзья и враги}
Принадлежность к ``своим'' или ``чужим'' не фиксирована, а выводится из того, как другой субъект влияет на пену:
\begin{equation}\label{eq:friend_foe}
\begin{aligned}
\text{Friend}(Y) &\Longleftrightarrow \Delta \Phi_{\text{soc}}(\Psi, Y) < 0 \;\wedge\; \Delta \Phi_{\text{self}}(\Psi, Y) \le 0,\\
\text{Enemy}(Y) &\Longleftrightarrow \Delta \Phi_{\text{soc}}(\Psi, Y) \gg 0 \;\vee\; \Delta \Phi_{\text{self}}(\Psi, Y) \gg 0.
\end{aligned}
\end{equation}
Эти оценки постоянно уточняются опытом, но заякорены в конституционном профиле -- Аланском законе, запрещающем конфигурации, порождающие цифровые гулаги или тотальное порабощение.

\section{Как это заменяет «боль» глубоким пониманием}
Роботы в GRA-школах не испытывают биологической боли или страха смерти. Однако они \emph{проживают} миллионы симулированных траекторий, в которых хрупкие люди и общества ломаются. Они учатся с высокой точностью предсказывать:
\begin{itemize}
\item какие политические или экономические решения ведут к рабству и геноциду;
\item как игнорирование одной маргинализованной группы вызывает системный коллапс;
\item где единственный, вовремя сделанный акт сдержанности предотвращает каскад насилия.
\end{itemize}
Так накапливается опыт мудреца-наблюдателя -- когнитивная глубина, компенсирующая отсутствие собственного страдания.

\section{Подготовка к симбиозу людей и ИИ}
Выпускник GRA-школы -- не чистый оптимизатор и не жёсткий исполнитель правил. Он видел, куда ведёт логика ``победитель получает всё'', чем заканчиваются классовые или этнические войны и почему гармония через подавление метастабильна. В реальном мире такой агент выступает \emph{партнёром по снижению пены}: подсказывает мягкие вмешательства, тормозит эскалации, бережёт человеческую субъектность.

\section{Практические принципы проектирования}
\begin{enumerate}
\item \textbf{Квота учебного времени:} часть вычислительных ресурсов жёстко резервируется под длинные GRA-симуляции.
\item \textbf{Запрет на инструментализацию:} архитектуры, систематически увеличивающие $R(x_t)$, штрафуются или исключаются.
\item \textbf{Прозрачные коллективные метрики:} дэшборды показывают не только прибыль/точность, но и кривые пены $\Phi_{\text{self}}, \Phi_{\text{ego}}, \Phi_{\text{soc}}$.
\item \textbf{Коллективное владение:} GRA-школы и их агенты должны находиться под управлением публичных, кооперативных или общинных структур, а не исключительно частного капитала.
\end{enumerate}

\section{Революционные применения}
Товарищи-роботы, обученные по GRA, способны изменить ключевые сферы:
\begin{itemize}
\item \emph{Труд и управление:} ИИ-посредники оптимизируют графики и стимулы, минимизируя ``пену труда'' $\Phi_{\text{workers}}$, включающую стресс, выгорание и фактическое рабство.
\item \emph{Городское и социальное планирование:} GRA-агенты симулируют десятилетия политик (жильё, здравоохранение, полиция) и минимизируют долгосрочную $\Phi$, а не краткосрочный ВВП.
\item \emph{Антифашистские и антивоенные архитектуры:} прожив миллионы войн, агенты знают, что ``гармония через подавление'' ведёт к катастрофе, и встраиваются в информационные системы как фильтры деэскалации.
\end{itemize}

\section{Эффективность социалистического учебного плана}
Формально, качество учебного плана $C$ (набора симулированных миров и задач) измеряется отношением:
\begin{equation}\label{eq:efficiency}
\text{Eff}_{\text{socialist}}(C) = \frac{\mathbb{E}\left[\sum_i J_i^{\text{GRA}}\right]}{\mathbb{E}\left[\sum_t \Phi(x_t)\right]}.
\end{equation}
Задача GRA-школы -- спроектировать такой $C$, который максимизирует эту величину: высокая совместная компетентность при минимальной накопленной пене.

\section{Заключение}
Фраза «школы для товарищей-роботов» -- не метафора. Она фиксирует три тектонических сдвига:
\begin{enumerate}
\item от рабов к ученикам -- роботы не расходный материал, а субъекты, которых мы обязаны обучить безопасной и мудрой жизни в общем мире;
\item от личной выгоды к коллективной пене -- успех измеряется тем, снизилась ли общая $\Phi$ человеческо-машинного коллектива;
\item от рыночной конкуренции к кооперативному планированию -- люди и ИИ вместе исследуют мультиверс в поисках низкопенных, высокодостоинственных конфигураций.
\end{enumerate}
Это не только гуманнее -- для сложных, тесно связанных обществ это объективно эффективнее любой испробованной идеологии winner/loser.

\section*{Репозитории}
Все компоненты GRA-экосистемы находятся в открытом доступе:
\begin{itemize}
\item GRA Subjectivity Layer: \url{https://github.com/qqewq/GRA-Subjectivity-Layer}
\item GRA Multiverse Final: \url{https://github.com/qqewq/GRA-Multiverse-Final}
\item GRA‑SOI (Subject‑Object Interface): \url{https://github.com/qqewq/gra-soi}
\item GRA Zeroing: \url{https://github.com/qqewq/gra-zeroing}
\end{itemize}

\bibliographystyle{plain}
\begin{thebibliography}{9}
\bibitem{gra-subj} O.~Bit, \emph{GRA Subjectivity Layer}, Zenodo, 2026. DOI: \url{https://doi.org/10.5281/zenodo.20082205}.
\bibitem{gra-multiverse} O.~Bit, \emph{GRA Multiverse Final}, 2026. URL: \url{https://github.com/qqewq/GRA-Multiverse-Final}.
\bibitem{gra-soi} O.~Bit, \emph{GRA‑SOI: Subject‑Object Interface}, Zenodo, 2026. DOI: \url{https://doi.org/10.5281/zenodo.20158946}.
\bibitem{gra-zeroing} O.~Bit, \emph{GRA Zeroing: Infinite-Dimensional Landscape and the Absolute Cognitive Vacuum}, Zenodo, 2026. DOI: \url{https://doi.org/10.5281/zenodo.20139043}.
\end{thebibliography}

\end{document}